\scanpage{55}
\mspage{6}
% archive foliation: 54
donner une section de $(X\setminus\Delta)/\underline V$ \deleted{\unsure{de}}, en prenant \inserted{pour $I$} l'image inverse de cette section dans $X\setminus\Delta$.

Mais la donnée d'une structure octaédrale équivaut à celle d'un \deleted{\unsure{$I$}} rev. étale de degré quatre de $S$ (à ne pas confondre avec le $\underline I$ de tantôt, appelons le $\underline K$), ou encore, à celle de

 a) un schéma en groupes $V$ sur $S$, loc. isom. à $(\mathbb D_2)_S$ -- ou encore un revêtement fini étale $\underline J\ (\unsure{\Leftarrow|=}V^*)$ de degré $3$ de $S$, ou \deleted{\unsure{de}} enfin, à celle d'un \deleted{\unsure{ce}} fibré en droites projectives $\bar\Sigma=\Sigma_{\underline J}$ sur $S$, muni d'un diviseur \inserted{relatif} étale $D_J$ de degré $3$ (\unsure{s. isom.} \deleted{\unsure{lui}} à $\underline J$)

 b) d'un torseur $\underline K$ sous $\underline V$, ou ce qui revient au $=$, d'un revêtement \deleted{galoisien} \inserted{\uline{principal}} $X^*$ de $\Sigma^*_{\underline J}=\Sigma_{\underline J}\setminus\deleted{\unsure{J}}D_{\underline J}$ de groupe \deleted{$\unsure{H^0(S,\mathscr H^1(\ ,V))}$} \inserted{$(\underline V_{\underline J})_{\Sigma^*_{\underline J}}$}, \deleted{géométriquement connexe, de degré \unsure{fibre}} \inserted{correspondant à $\unsure{\xi}$ un élément de
\[
H^0(S,\mathscr H^1(\Sigma^*_{\underline J}/S,\underline V))\simeq\operatorname{Hom}(\underline V_{\underline J},\underline V)\deleted{\unsure{\simeq\mathfrak S_3}}
\]
qui est l'identité $\underline V_{\underline J}\longrightarrow\underline V=\underline V_{\underline J}$} -- la correspondance entre $\underline K$ et $X$ étant donnée par la formule
\[
(11)\qquad X^*=\deleted{\unsure{X_{\underline J,0}\times}}\quad \underline K\wedge_V\unsure{\widetilde\Sigma}_{\underline J,2}^*,
\]
où $\unsure{\widetilde\Sigma}_{\underline J,2}^*$ est le revêtement principal dé

\margin{\uline{NB} Le grpt. \unsure{local des} $(\pi_1)_{\mathrm{ab}/2}$ de $\Sigma^*_{\underline J}$ sur $S$ est can. isom. à $\underline V_{\underline J}$\ldots}
